\section{Metodología}

La gestión se basa en un enfoque adaptativo. Esta denominación describe la forma de gestionar el proyecto y no un tipo de modelo predictivo. Se elige porque los requisitos, la calidad y representatividad de los datos, el desempeño de las alternativas, los riesgos y las necesidades de validación pueden evolucionar durante el proyecto.

En lugar de asumir una planificación completamente estable, el enfoque adaptativo organiza ciclos incrementales de planificación, análisis, evaluación y ajuste. Cada ciclo mantiene entregables, responsables e hitos definidos, pero permite incorporar evidencia antes de decidir las etapas siguientes.

\subsection{Ciclos de gestión propuestos}

\begin{description}[style=nextline]
  \item[Ciclo 1] Comprensión del problema, objetivos, actores y criterios de éxito. Hito: alcance inicial aprobado.
  \item[Ciclo 2] Comprensión y evaluación de datos: calidad, representatividad y privacidad. Hito: factibilidad inicial de datos validada.
  \item[Ciclo 3] Diseño y evaluación de Cloud, Local e Híbrida, herramientas, costos preliminares y riesgos. Hito: alternativas técnicamente viables.
  \item[Ciclo 4] Evaluación integral técnica, ética, financiera, legal y de riesgos. Hito: alternativa preferida seleccionada.
  \item[Ciclo 5] Planificación de un eventual piloto: recursos, seguridad, operación y mantenimiento. Hito: plan de ejecución aprobado.
  \item[Ciclo 6] Evaluación, retroalimentación, ajustes y lecciones aprendidas. Hito: decisión de continuidad o escalamiento.
\end{description}

La planificación debe incluir revisiones, hitos y capacidad de ajuste; las finanzas deben incorporar horas de revisión e iteración; los riesgos deben reevaluarse durante los ciclos; y la ética debe considerar sesgos y explicabilidad de forma continua, no solo al cierre.
